\subsubsection{Zero-Knowledge}
\label{appendix:const-zk}

\begin{theorem} \label{thm:zero-knowledge}
  The generic construction is zero-knowledge in the honest majority setting
  ($\tcor < n/2$) in the random oracle model,
  against a PPT adversary $\adv$ controlling up to $t-1$ participants,
  assuming the generic construction itself is weakly robust,
  the \AVSS $\avss$ achieves aggregated secrecy,
  and the \ENIKE $\enike$ is unrecoverable.
%  \ian{Thm 1 specified an honest majority.  Is that not a requirement here as well?}
  % CK- it says that in the above statement

  Concretely, for every adversary $\adv$ that attacks the construction,
  there exist adversaries $(\bdv_1, \bdv_3, \bdv_4)$ that run in about the same time as $\adv$,
  and algorithms $\simkeygen_i, i \in \{0,\ldots,2 \}, \simfin$,
  such that
  \begin{align} \label{eq:zk-thm}
  \begin{split}
      \advantzk{\dkg,\adv,\simkeygen_i,\simfin}(\lambda) & \leq \advantsecrecy{\avss[\Hash_1],\bdv_1}(\lambda)\ +
        \advant{rec}{\enike,\bdv_3}(\lambda) \cdot q_h \\
      & \quad + \advantaggsecrecy{\avss[\Hash_1], \bdv_4, \simshare,\simgettweak}(\lambda)
      \end{split}
  \end{align}
  where $\simshare, \simgettweak$ are as defined for aggregated secrecy of the \AVSS,
  as in Definition~\ref{defn:aggsecrecy},
  and $q_h$ denotes the number of queries $\adv$ is allowed to make to $\Hash_1$.
\end{theorem}


\begin{proof}
  We now complete the proof of Theorem~\ref{thm:zero-knowledge} via a sequence of games.


\medskip

\noindent \textit{\underline{Game 0.}}
This is the DKG zero-knowledge game when $b=0$ as in Figure~\ref{fg:zk},
applied to the construction in Figure~\ref{fg:generic-perfect}.

Let $W_0$ be the event that $\adv$ outputs $1$ in Game~0. Then,
\[ \advant{zk}{\dkg,\adv,\simkeygen_i,\simfin}(\lambda)  = \Pr[W_0]    \]

Let $\cdv$ be the algorithm that simulates the DKG zero-knowledge game to $\adv$.
We now describe in further detail how $\cdv$ simulates the game.

\medskip

\noindent \textit{Setup.}
To begin, $\cdv$ initializes $Q_1 \gets \emptyset$ to simulate $\Hash_1$
and $Q_2 \gets \emptyset$ to simulate $\Hash_2$,
in addition to initializing each honest party as shown in Figure~\ref{fg:zk}.
$\cdv$ handles $\adv$'s random oracle queries by lazy sampling, as follows:

  \begin{itemize}
    \item[$\underline{\Hash_1}$]: When the adversary queries $\Hash_1$ on inputs $(\pubaggregateset_i, \auxstring_i)$,
      the environment checks if it is in $Q_1$, and if so, returns the corresponding $\tweak_i$.
      Otherwise, the environment randomly samples $\tweak_i$ from its respective domain,
      sets $Q_1[(\pubaggregateset_i, \auxstring_i)] = \tweak_i$, and then returns $\tweak_i$.
    \item[$\underline{\Hash_2}$]: When the adversary queries $\Hash_2$ on inputs $\{\blinding_1,\ldots,\blinding_n \}$,
      the environment checks if it is in $Q_2$, and if so, returns the corresponding $\auxstring_i$.
      Otherwise, the environment randomly samples $\auxstring_i$ from its respective domain,
      sets $Q_2[\{\blinding_1,\ldots,\blinding_n \}] = \auxstring_i$, and then returns $\auxstring_i$.
  \end{itemize}

$\cdv$ performs all signing oracle queries honestly.

\medskip

\noindent \textit{\underline{Game 1.}}
The only difference between Game~0 and Game~1 is that in Game~1,
$\cdv$ performs $\simkeygen_0, \ldots, \simkeygen_{\numrounds-1},
\simfin$ for participant $\simid \rgets \honest$ with respect to the challenge $\dkgchal$.
More specifically,
when performing $\simkeygen_0$,
$\cdv$ employs $(\{ \share_{\tau,j}^* \}_{j \in \corrupt},\allowbreak\vsscommit_\simid^*) \gets \avss[\Hash_1].\simshare(\lambda,n,t,\simid,\dkgchal)$ for an honest participant $\simid$ sampled from the set of honest participants at random,
such that Equation~\ref{eq:g2-sim} holds:
  \begin{equation} \label{eq:g2-sim}
    \vess.\derivepshare(0, \vsscommit_\simid^*) = \dkgchal
  \end{equation}

$\cdv$ performs $\simkeygen_1, \simkeygen_2,\simfin$ as in the real protocol.
However, one consequence of $\cdv$ simulating secret sharing with respect to $\dkgchal$ is that
$\cdv$ cannot correctly derive each honest players' $\SK_i$, for all $i \in
\honest$ when performing $\simfin$.
This however does not impact the ability of $\cdv$ to perfectly
simulate all honest players, because $\cdv$ is never required to perform
operations with respect to $\dkgchal$ directly.

If the \AVSS is secret,
then Game~0 and Game~1 are indistinguishable to $\adv$.

\medskip

\noindent \textit{Reduction to \AVSS (Non-Aggregated) Secrecy.}
We will construct a reduction $\bdv_1$ which is an adversary against the \AVSS (non-aggregated) secrecy game as defined in Definition~\ref{defn:vss-secrecy},
such that $\bdv_1$ simulates either Game~0 or Game~1 to $\adv$.
When the \AVSS hidden bit $b=0$, $\bdv_1$ ends up simulating Game~0;
when the \AVSS hidden bit $b=1$, $\bdv_1$ ends up simulating Game~1.

$\bdv_1$ is constructed similar to Game~0, with the following exception.
$\bdv_1$ simply performs $\dkgkeygen_1$ for each honest non-simulated party $i \in \honest, i \neq \simid$.
However, for the simulated honest party $\tau$,
instead of following the protocol honestly,
$\bdv_1$ instead employs the challenge received as input $(\{(j, \share_j)\}_{j \in \corrupt}, \vsscommit) $
from the \AVSS secrecy experiment.
In particular, it sets $\vsscommit_\simid \gets \vsscommit$,
and $(\{(j, \share_{\simid,j} )\}_{j \in \corrupt}\allowbreak\gets (\{(j, \share_j)\}_{j \in \corrupt}$.
It then follows the remainder of the protocol honestly.
While it will not be able to derive the honest player's secret keys in $\dkgkeygen_3$,
this will not have an impact on its ability to correctly simulate.

%\douglas{This reduction is very thin -- too much of a skevtch in my opinion.  You say that the reduction ``queries its \AVSS secrecy oracle'', but in fact there are three oracles it could be querying -- Oget, Orev, Otweak. Which one?  How does it use the other oracles?}
% CK- this is a reduction to the non-aggregated game, which we can do because
% the AVSS is an extension of the VSS scheme.
% In that game,

The simulation of Game~0 by $\bdv_1$ is identical to when $b=0$ in the \AVSS secrecy game.
Hence, the only distinction between Game~0 and Game~1 is when $b=1$ in the \AVSS secrecy game.

When $\adv$ outputs $b'$ as its guess,
then $\bdv_1$ outputs $b'$.
Hence, when $\adv$ distinguishes Game~0 from Game~1,
$\bdv_1$ wins the \AVSS secrecy game.




\medskip

\noindent \textit{Difference between Game~0 and Game~1.}
If the \AVSS is secret, then
the difference between Game~0 and Game~1 is negligible.
Let $W_1$ be the event that $\adv$ wins in Game~1.
Then
\begin{equation} \label{eq:game1-zk}
  \bigabs{ \Pr[W_1]  - \Pr[W_0] } \leq   \advant{sec}{\avss,\bdv_1}(\lambda)
\end{equation}


\medskip

\begin{figure}[p]
\begin{pchstack}[boxed,center,space=1em]
\pcvspace
\procedure[]{$\dkgkeygen_1(\party_i, \invecp, \invecb)$} {
  \text{Parse inputs honestly as defined in the protocol} \\
  \pcfor j \in \set{n}, j \neq i \pcdo \\
  \pcind \pcif \avss[\Hash].\verify(i, \share_{ji}, \vsscommit_j) \neq 1 \\
  \pcind \textbf{or }\pcif \avss[\Hash].\derivepshare(0, \vsscommit_j) \neq \pk_j \\
  \pcind \pcind \party_i.\status = \abort \\
  \pcind \pcind \pcreturn (\party_i, \bot, \fail) \\
  \party_i \gets \party_i \cup \{ (i, \share_{ji}) \}_{j \in \set{n}} \\
  %
  \framebox{$\recoveryset_k = \{(j,\share_{kj}) \}_{k \in \corrupt, j \in \honest}$} \\
  \framebox{$\secretone_k \gets \avss.\recover[\Hash](t, \recoveryset_k, \vsscommit_k) $} \\
  \framebox{$\pcreturn \bot\ \pcif \secretone_k = \fail$} \\
  \framebox{$\blinding_k \gets \enike.\getsharedkey[\Hash](\secretone_k, \commitment_k), \forall\ i \in \corrupt $} \\
  \framebox{$\pubaggregateset =  \{\vsscommit_j\}_{j \in \set{n}} $} \\
  \framebox{$\auxstring \gets \{\blinding_j \}_{j \in \set{n}}$} \\
  \pcreturn (\party_i,\ \bot,\ \accept)
}
\end{pchstack}
    \caption{Game 2 for the Zero Knowledge Proof for Generic Construction.
    }
  \label{fg:zk-g3}
\end{figure}

\noindent \textit{\underline{Game 2.}}
The difference between Game~1 and Game~2 is code movement from
$\finalize$ to $\dkgkeygen_1$,
as shown in Figure~\ref{fg:zk-g3},
which shows code movement.
In particular, the code that is moved is highlighted,
moved from $\finalize$ to $\dkgkeygen_1$.

Rather than learning corrupted players' blinding factors as input to $\finalize$,
$\cdv$ instead derives blinding factors for all corrupted parties in $\dkgkeygen_1$,
(importantly, before it reveals the honest players' blinding factors to $\adv$)
using the shares from all honest players.
If $\cdv$ is unable to recover these blinding factors,
then it returns $\bot$.

However, because $\cdv$ simulates at minimum $t$ honest players (honest majority model),
it can perform $\avss.\recover$ without the involvement of corrupted parties.


\medskip

\noindent \textit{Difference between Game~1 and Game~2.}
Let $W_2$ be the event that $\adv$ wins in Game~2. Then,

\begin{equation} \label{eq:game3-zk}
  \bigabs{ \Pr[W_2]} = \bigabs{ \Pr[W_1]  }
\end{equation}




\medskip

\noindent \textit{\underline{Game 3.}}
The only difference between Game~3 and Game~2 is that if the adversary queries
the random oracle $\Hash_1$ with respect to any of the honest parties' blinding values $\{\blinding_i \}_{i \in \honest}$ before the adversary is provided with the witness values $\{ \witness_i\}_{i \in \honest}$ as input to $\dkgkeygen_3$,
then the game aborts.

If the \ENIKE is session-key unrecoverable,
the additional advantage to $\adv$ is negligible.

\medskip

\noindent \textit{Reduction to \ENIKE Unrecoverability.}
We will construct a reduction $\bdv_3$ which is an adversary against the \ENIKE
unrecoverability game, such that $\bdv_3$ simulates Game~3 to $\adv$.

$\bdv_3$ begins by receiving $(\PK_1,\PK_2)$ in the \ENIKE unrecoverability game.
It then simulates Game~3 to $\adv$ in exactly the same manner as in Game~2,
with the following exception.
$\bdv_3$ simply performs $\dkgkeygen_0$ honestly for each honest party $i \in \honest, i \neq \simid$.
However, for honest party
$\simid$, $\bdv_3$ instead simulates the protocol via $\simkeygen$;
it sets $\PK_\simid \gets \PK_1$ and $\commitment_\simid \gets \PK_2$
$\bdv_3$ then simulates secret sharing via $\simshare$ with respect to $\PK_\simid$.
% CK - this is why we need the hop to VSS secrecy
From Game~1, we know that doing so is indistinguishable to $\adv$.
$\bdv_3$ then follows the remainder of $\dkgkeygen_1$ as for all the other honest parties.

When simulating $\oraclekg_2$,
$\bdv_3$ acts honestly for all honest parties $i \in \honest, i \neq \simid$.
However, for participant $\simid$, $\bdv_4$ does the following.
Because it does not know the corresponding $\witness_c$ to $\commitment_c$,
it instead must guess.
It looks at all of the queries made to $\Hash_1$ that $\adv$ has made up to that point,
and randomly selects one of the inputs.
Let this randomly selected input be $(\{\blinding_j \}_{j \in \set{n}}^*, \auxstring^*)$.

$\bdv_3$ then picks the simulated party's element $\blinding_\simid \gets \{\blinding_j \}_{j \in \set{n}}^*$,
and outputs $\sharedkey' = \blinding_\simid$ as its guess for the \ENIKE unrecoverability game.
%
If $\adv$ queries $\Hash_1$ with $\blinding_\simid$ before $\bdv_3$ has
completed $\oraclekg_2$,
then $\bdv_3$ wins the \ENIKE game with probability $\frac{1}{\lvert \Hash_1 \rvert}$,
where $\lvert \Hash_1 \rvert$ denotes the order of the co-domain of $\Hash_1$.

\medskip

\noindent \textit{Difference between Game~2 and Game~3.}
Let $W_3$ be the event that $\adv$ wins in Game~3.
The additional ability for $\adv$ to win Game~3 is then bounded by its ability to win the \ENIKE unrecoverability game.
Hence,

\begin{equation} \label{eq:game3-secrecy}
  \bigabs{\Pr[W_3] - \Pr[W_2] } \leq
\advant{rec}{\enike,\bdv_3}(\lambda) \cdot q_h
\end{equation}

\medskip

\noindent \textit{\underline{Game 4.}}
The only difference between Game~3 and Game~4 is that in Game~4,
$\cdv$ simulates secret sharing for participant $\simid \rgets \honest$ with respect to the challenge $\dkgchal$.
From Game~1,
we know that $\cdv$ can do so in such a way that is indistinguishable to $\adv$.
However, in Game~4,
it additionally does so in such a way that Equation~\ref{eq:program-tweak} holds after performing $\finalize$.
\begin{equation} \label{eq:program-tweak}
  \avsscommit[\Hash_1].\derivepshare(0, \avsscommit) = \dkgchal
\end{equation}


More specifically,
in $\oraclekg_0$,
$\cdv$ performs
\[ (\{ (j, \share_{\simid j}) \}_{j \in \corrupt}, \vsscommit_\simid) \rgets \avss[\Hash_1].\simshare(\lambda, \corrupt, \dkgchal)\,. \]
It then follows the remainder of $\dkgkeygen_0$ as in Game~3.

Then,
in $\oraclekg_3$,
it derives $\tweak'$ as
$\tweak' \rgets \avss[\Hash_1].\simgettweak(\pubaggregateset, \party_\simid)$.
It then programs $\Hash_2(\pubaggregateset, \auxstring)$ with $\tweak'$.
From Game~3, we know that $\cdv$ can derive $\auxstring$ in $\oraclekg_3$,
but $\adv$ has only negligible probability of doing so.


We will see that, if the \AVSS is securely aggregatable,
then the additional advantage to the adversary in Game~4 is negligible.

\medskip

\noindent \textit{Reduction to \AVSS Aggregated Secrecy.}
We will construct a reduction $\bdv_4$ which is an adversary against the AVSS aggregated secrecy game,
such that $\bdv_4$ simulates either Game~3 or Game~4 to $\adv$.
When the hidden bit $b$ in the AVSS aggregated secrecy game is $b=0$, $\bdv_4$ ends up simulating Game~3;
when $b=1$,  $\bdv_4$ ends up simulating Game~4.

$\bdv_4$ is constructed similar to Game~3, except as follows.
When $\adv$ queries $\oraclekg_1$,
for each honest player $k \in \honest$,
$\bdv_4$ queries the AVSS oracle $\getshare()$,
receiving $(\{ \share_j \}_{j \in \honest}, \allowbreak \vsscommit^*)$ in return.
$\bdv_4$ then uses these outputs to simulate participant $\simid$ for $\oraclekg_1$.
%
$\bdv_4$ follows the protocol honestly for all other players,
and  submits all $(\{ \share_{kj} \}_{j \in \honest}, \vsscommit_k), k \in \set{n} \setminus \simid$
to the AVSS oracle $\recvshare$.

Then, $\bdv_4$ receives $\auxstring'$,
and then queries the \AVSS aggregated secrecy oracle $\Ogettweak$ on inputs $(\pubaggregateset, \auxstring')$,
 receiving $\tweak'$ in return.
$\bdv_4$ then programs its own $\Hash_1(\pubaggregateset, \auxstring) = \tweak'$.
From Game~3,
$\adv$ can query for its output with negligible probability before $\bdv_4$ programs it.
$\bdv_4$ then follows the remainder of Game~3 as before.

When $\adv$ outputs $b'$,
then $\bdv_4$ outputs $b'$.
Hence, when $\adv$ distinguishes Game~3 from Game~4,
$\bdv_4$ wins the \AVSS aggregated secrecy game.


\medskip

\noindent \textit{Difference between Game~3 and Game~4.}
Let $W_4$ be the event that $\adv$ wins in Game~4.
The difference between $W_3$ and $W_4$ is bounded by the advantage of $\adv$ in the \AVSS aggregated security game.
\begin{equation} \label{eq:game4-zk}
  \bigabs{\Pr[W_4] - \Pr[W_3] } \leq
  \advant{asec}{\avss[\Hash_1], \bdv_4, \simshare,\simgettweak}(\lambda)
\end{equation}


\medskip

\noindent \textit{Finishing the Proof.}
The resulting public key in Game~4 is identical to the challenge $\dkgchal$,
and hence is identical to the DKG zero-knowledge game when $b=1$.
The adversary then has at most negligible advantage in winning the DKG zero-knowledge game.
Concretely, combining the advantage of the adversary in (\ref{eq:game1-zk})-(\ref{eq:game4-zk}) proves Equation~\ref{eq:zk-thm}.
This concludes the proof.
\end{proof}

